Sei $M$ die Masse von Last und Karre, $a$ der Abstand ihres Schwerpunkts und $b$ der Abstand der Griffe von der Radachse.

\begin{abcliste}
\abc
Nimm die Radachse als Drehachse: Die Kraft des Rads hat dann keinen Hebelarm. Last und Hände liegen auf derselben Seite der Achse (ein einseitiger Hebel), drehen die Karre aber in entgegengesetzte Richtungen:
\begin{align}
 F\cdot b &= M\,g\cdot a \\
 F &= \FHaF \\
   &= \frac{\ML\cdot\g\cdot\aL}{\bG} \\
   &= \FHa \approx \result{\FHaP{0}{2}}
\end{align}

\abc
Die Kräfte heben sich auf: Hände und Rad tragen zusammen die Gewichtskraft.
\begin{align}
 F_\mathrm{R} &= \FRF \\
   &= \ML\cdot\g-\FHa \\
   &= \FR \approx \result{\FRP{0}{2}}
\end{align}
Das Rad trägt den grössten Teil der Last.
\end{abcliste}
